\section{Discussion}
\label{sec:discussion}

The Fisher--Rao bounds proved here are not summarized by one condition number.
Covariance underdispersion, global nonquadratic localization, and local
mean--covariance coupling are three different effects.  Their separation
explains several facts that otherwise look contradictory.  Fisher--Rao is
affine invariant, yet its global rate can depend on the initial covariance.
A transport warm start can eliminate that dependence, yet a
\(\kappastar\)-scale lower bound remains.  Optimizer whitening can remove
extrinsic anisotropy, yet the local score generator can still be slow because
the whitened target is non-Gaussian.

Three sharpness statements should be read with their scopes attached.  The
spiral gives a continuous-time \(\Omega(\kappastar)\) localization time in
dimension at least two for a specified anisotropic initialization class.  The
bump train gives a one-dimensional \(\Omega(\kappa/h)\) lower bound for both
fixed-step maps, with \(h=\gamma/\kappa\) and
\(0<\gamma\leq1/2\), for a constant-factor reduction after covariance
burn-in.  The convex ridge gives
\(\gstar=\Theta(\kappastar^{-1/2})\) with
\(d=\Theta(\kappastar^2)\).  None of these statements alone proves a
fixed-dimensional lower bound with the same dependence, and the ridge does not
saturate the \(\sqrt d\log\kappastar\) branch of the universal estimate.

The stochastic results carry equally important qualifications.  A constant-
step theorem gives geometric convergence to an explicit residual upper bound;
it neither identifies the stationary error nor, when that residual is nonzero,
certifies arbitrary accuracy at fixed step and batch size.  Under
Hessian-Lipschitzness, the decreasing-step result instead certifies decay to
zero at an \(O(1/N)\) rate.  Locally,
contraction is first proved for the process killed at the boundary of a
certified region.  A separate finite-horizon argument turns this into a
high-probability non-exit statement, at the price of deeper entry and a batch
size that grows with the requested horizon under the present proof.  These
distinctions are visible in the theorem statements rather than being deferred
to caveats.

The classification of affine-invariant metrics explains why Fisher--Rao is a
natural default, not why it is fastest on every target.  For \(d\geq2\), the
parameters \((\eta,\omega,\tau)\) set independent time scales for the mean,
traceless covariance, and trace covariance modes on a Gaussian target.
Fisher--Rao is the unique balanced choice up to scale.  A deliberately
unbalanced metric can accelerate a known dominant mode, as the trace experiment
illustrates.  A complete sharp global and stochastic theory for every member of
this family would require new arguments; it is not implied by the Gaussian
modal calculation.

The non-log-concave example marks a structural boundary.  Negative curvature
is multiplied by the covariance in the Fisher--Rao dynamics.  A uniform Hessian
bound therefore does not prevent an unprojected covariance cascade or a pole in
the KL resolvent.  Spectral clipping changes the problem to constrained
optimization and restores a useful \(O(1/N)\) displacement guarantee, but the
displacement can vanish at an active constraint while the unconstrained
gradient remains large.

Several problems remain open.  The most immediate is the optimal local spectral
rate in fixed dimension.  A second is a closed-form general-target radius for
the discrete maps, replacing the compact-hull second-derivative supremum.
Sharper stochastic local stepsizes and an infinite-horizon stability mechanism
that does not require horizon-dependent batching would make the local theory
more practical.  Finally, general affine-invariant preconditioning and a full
Wasserstein--Fisher--Rao theory are natural continuations, but they should be
developed without obscuring the scoped sharp Fisher--Rao bounds established
here.
